\documentclass[11pt,a4paper]{article}

\usepackage[margin=2.2cm]{geometry}
\usepackage[T1]{fontenc}
\usepackage[utf8]{inputenc}
\usepackage{amsmath}
\usepackage{amssymb}
\usepackage{array}
\usepackage{booktabs}
\usepackage{hyperref}
\usepackage{xcolor}

\newcommand{\done}{\textbf{\textcolor{green!50!black}{[x]}}}
\newcommand{\wip}{\textbf{\textcolor{orange!80!black}{[~]}}}
\newcommand{\todo}{\textbf{\textcolor{gray}{[ ]}}}
\newcommand{\code}[1]{\texttt{#1}}

\title{HTW Dynamic Maze: strategie JEPA, A* distillation et bootstrap WM}
\author{HackTheWorld(s) -- EB-JEPA}
\date{\today}

\begin{document}
\maketitle

\section{Objectif}

Le pivot \emph{Dynamic Maze} transforme le benchmark maze statique en probleme
partiellement observe et stochastique. Le labyrinthe contient des portes qui
s'ouvrent et se ferment aleatoirement; l'agent ne voit qu'une fenetre locale
(\emph{fog-of-war}). Le but est de montrer qu'un world model (WM) peut exploiter
la dynamique apprise du monde, plutot que seulement optimiser une distance brute
dans l'espace latent.

Deux familles de methodes sont suivies:

\begin{enumerate}
  \item une baseline forte mais moins pure: entrainer le JEPA sur des trajectoires
        A* replannifiees, puis distiller des sous-buts \(N=2\) et \(N=4\);
  \item une version plus ``world model'': collecter des trajectoires par
        exploration locale sans A*, puis apprendre une value/reachability head par
        hindsight TD.
\end{enumerate}

\section{Structure du monde}

Chaque episode est construit ainsi:

\begin{enumerate}
  \item generer un labyrinthe DFS \(21\times 21\);
  \item choisir des cellules-murs candidates et les transformer en portes
        stochastiques;
  \item initialiser chaque porte ouverte avec probabilite \(p_0=0.35\);
  \item a chaque action, chaque porte toggle avec probabilite \(q=0.04\);
  \item rendre une observation locale sous fog-of-war.
\end{enumerate}

L'observation a quatre canaux:

\begin{center}
\begin{tabular}{clp{0.62\linewidth}}
\toprule
Canal & Nom & Sens \\
\midrule
0 & agent & point gaussien localisant l'agent. \\
1 & murs visibles & murs/portes fermees visibles sous fog. \\
2 & inconnu & masque des cellules cachees par le fog-of-war. \\
3 & portes visibles & cellules identifiees comme portes quand elles sont visibles. \\
\bottomrule
\end{tabular}
\end{center}

La probabilite marginale qu'une porte soit ouverte apres \(t\) deplacements est:

\[
  P_t(\mathrm{open}) =
  \frac{1}{2} + \left(p_0-\frac{1}{2}\right)(1-2q)^t.
\]

Avec \(p_0=0.35\) et \(q=0.04\), la probabilite remonte vers \(0.5\): environ
\(46.4\%\) apres 17 deplacements et \(49.9\%\) apres 65 deplacements.

\section{Donnees et tenseurs}

Le loader renvoie des fenetres de trajectoires:

\begin{center}
\begin{tabular}{p{0.28\linewidth}p{0.62\linewidth}}
\toprule
Tenseur & Sens \\
\midrule
\(\mathbf{x}\in\mathbb{R}^{B\times 4\times T\times 63\times 63}\)
& sequence d'observations dynamic-maze normalisees. \\
\(\mathbf{a}\in\mathbb{R}^{B\times 2\times T}\)
& actions cardinales en coordonnees pixel \((dr,dc)\). \\
\(\mathbf{p}\in\mathbb{R}^{B\times 2\times T}\)
& positions normalisees de l'agent, utilisees pour le probe et certains labels. \\
\code{wall\_x}, \code{door\_y}
& champs dummy pour garder la signature commune du repo. \\
\bottomrule
\end{tabular}
\end{center}

\section{Pseudo-code: generation A*}

La premiere version exploite A* comme professeur de donnees. A* est recalcule a
chaque pas sur la grille courante, donc il voit les portes ouvertes/fermees au
moment present. Il ne predit pas les toggles futurs.

\begin{quote}
\small
\begin{verbatim}
function SAMPLE_ASTAR_TRAJECTORY(config, rng):
    maze = DFS_RANDOM_MAZE(config.maze_height, config.maze_width)
    gates = SAMPLE_GATE_CELLS(maze, config.num_doors)
    gate_state[i] ~ Bernoulli(config.door_initial_open_prob)
    agent = start_cell
    goal  = goal_cell

    states, actions, locations = [], [], []

    for t in 0 .. config.n_steps - 1:
        obs_t = RENDER_FOG_OBSERVATION(maze, gates, gate_state, agent, goal)
        grid_t = APPLY_OPEN_GATES(maze, gates, gate_state)

        # A* professor: shortest path on the current full grid.
        path_t = ASTAR(grid_t, agent, goal)
        if path_t exists:
            action_t = FIRST_CARDINAL_STEP(path_t)
        else:
            action_t = ZERO_ACTION

        states.append(obs_t)
        actions.append(action_t)
        locations.append(PIXEL_POSITION(agent))

        agent = STEP_ENV(agent, action_t, grid_t)
        for each gate i:
            with probability q: gate_state[i] = not gate_state[i]

    start = RANDOM_SLICE_START(n_steps, sample_length)
    return WINDOW(states, actions, locations, start, sample_length)
\end{verbatim}
\end{quote}

Dans le code, cette variante correspond a:

\begin{center}
\begin{tabular}{ll}
\toprule
Config & Valeur \\
\midrule
\code{teacher\_policy} & \code{oracle\_replan} \\
\code{layout\_use\_astar\_filter} & \code{true} \\
\code{sample\_length} & 17 \\
\code{n\_steps} & 65 ou 129 selon le run \\
\bottomrule
\end{tabular}
\end{center}

\section{Pseudo-code: entrainement JEPA A*}

Le JEPA apprend la dynamique action-conditionnee dans l'espace latent. L'encodeur
produit \(z_t=f_\theta(x_t)\), le predictif recurrent produit
\(\hat z_{t+k}=g_\phi(z_t,a_{t:t+k-1})\), puis on compare aux latents reels.

\begin{quote}
\small
\begin{verbatim}
function TRAIN_JEPA_ASTAR(config):
    loader = STREAM_DATASET(policy = oracle_replan)
    encoder   = ImpalaEncoder(input_channels = 4)
    predictor = RNNPredictor(action_dim = 2)
    regularizer = VICReg + temporal similarity + inverse dynamics
    probe_xy = MLP(z -> xy)

    for epoch in 1 .. E:
        for batch (x, a, loc) in loader:
            z_real = encoder(x)
            z_pred = predictor.rollout(z_real[:, t0], actions = a)

            L_pred = MSE(z_pred, stopgrad(z_real_future))
            L_reg  = cov/std/sim/idm regularization
            L_probe = MSE(probe_xy(stopgrad(z_real_t0)), loc_t0)

            L_total = L_pred + L_reg
            update(encoder, predictor, regularizer) on L_total
            update(probe_xy) on L_probe

    save encoder, predictor, probe_xy
\end{verbatim}
\end{quote}

La loss principale est:

\[
  \mathcal{L}_{pred}
  = \frac{1}{K}\sum_{k=1}^{K}
    \left\|\hat z_{t+k}-\mathrm{sg}(z_{t+k})\right\|_2^2.
\]

Avec les regularisations:

\[
  \mathcal{L}_{WM}
  = \mathcal{L}_{pred}
  + \lambda_{std}\mathcal{L}_{std}
  + \lambda_{cov}\mathcal{L}_{cov}
  + \lambda_{sim}\mathcal{L}_{sim}
  + \lambda_{idm}\mathcal{L}_{idm}.
\]

\section{Sous-buts A* distilles}

Apres entrainement du WM, on gele le JEPA et on entraine un MLP de sous-but. Pour
un instant \(t\), le label est la position \(N\) pas plus loin dans la trajectoire
observee:

\[
  y_t = p_{\min(t+N,T)}, \qquad
  \mathcal{L}_{sg}
  = \left\|h_\psi(z_t,p_T)-y_t\right\|_2^2.
\]

Pseudo-code:

\begin{quote}
\small
\begin{verbatim}
function TRAIN_SUBGOAL_HEAD(frozen_jepa, N):
    for batch (x, a, loc) in A*_trajectory_loader:
        z = frozen_jepa.encode(x)
        goal_xy = loc[:, :, -1]
        label_t = loc[:, :, min(t+N, T-1)]
        pred_t = SubgoalPredictor(z_t, goal_xy)
        update head on MSE(pred_t, label_t)
\end{verbatim}
\end{quote}

A l'evaluation, A* n'est plus appele. Le head predit un waypoint, puis le WM
teste les actions cardinales par rollout latent et choisit celle qui se rapproche
le mieux du waypoint.

\section{Resultats actuels}

Run starter:

\begin{quote}
\code{$EBJEPA_WORK/runs/dynamic\_maze\_jepa\_subgoal\_starter\_20260620\_034609}
\end{quote}

\begin{center}
\begin{tabular}{lrrrr}
\toprule
Methode & Success & Efficiency & Steps moyens & Blocked moves \\
\midrule
\code{fog\_optimistic} & 14.1\% & 0.134 & 437.7 & 0.0 \\
\code{memory\_optimistic} & 100.0\% & 0.909 & 79.7 & 0.0 \\
\code{oracle\_replan} & 100.0\% & 0.966 & 70.7 & 0.0 \\
JEPA subgoal \(N=2\) & 57.8\% & 0.403 & 280.8 & 103.3 \\
JEPA subgoal \(N=4\) & 73.4\% & 0.494 & 230.7 & 93.0 \\
\bottomrule
\end{tabular}
\end{center}

Lecture: le WM distille A* bat largement les baselines fog-only, mais reste loin
des solveurs symboliques optimistes. Le principal defaut est l'efficacite: il
atteint souvent le but, mais tente encore trop de mouvements bloques.

\section{Version pure WM: bootstrap sans A*}

Pour mieux coller au theme \emph{world models}, la prochaine variante supprime
A* de la collecte de donnees et du budget d'evaluation.

\paragraph{Collecte.}
On utilise \code{local\_goal\_frontier}: une politique myope qui ne regarde que la
carte visible, les cellules deja visitees et une faible heuristique Manhattan vers
le goal. Elle ne lance jamais de graph search.

\begin{quote}
\small
\begin{verbatim}
function LOCAL_GOAL_FRONTIER_POLICY(obs, memory):
    for each cardinal action a:
        if local next cell is blocked:
            score[a] = -infinity
        else:
            new_visible = fog_window(next_cell) - known_cells
            revisit = visit_count[next_cell]
            goal_bias = manhattan(next_cell, goal_cell)
            score[a] = novelty_bonus * |new_visible|
                       - revisit_penalty * revisit
                       - goal_bias_weight * goal_bias
                       + small_noise
    with probability epsilon:
        return random_cardinal_action()
    return argmax_a score[a]
\end{verbatim}
\end{quote}

\paragraph{Value hindsight.}
Dans une fenetre de trajectoire exploree, tout etat futur devient un pseudo-but.
On entraine une value head \(V_\eta(z_t,z_g)\) a predire le retour discount vers
ce pseudo-but, et on la calibre aussi sur les rollouts imagines du WM.

\begin{quote}
\small
\begin{verbatim}
function TRAIN_BOOTSTRAP_WM_VALUE(config):
    loader = STREAM_DATASET(policy = local_goal_frontier,
                            layout_use_astar_filter = false)
    train JEPA dynamics as before

    for batch (x, a, loc) in loader:
        z_real = jepa.encode(x)
        z_roll = jepa.rollout(x_0, a)

        for horizon h in {1,2,4,8,16,32}:
            goal = z_real[:, t+h]
            done = 1 if h == 1 else 0
            target = done + gamma * (1-done) * V_target(z_next, goal)

            L_value += MSE(V(z_real_t, goal), target)
            L_value += MSE(V(z_roll_t, goal), target)

        update value head
        EMA-update target value head
\end{verbatim}
\end{quote}

\paragraph{Evaluation.}
Le controle est A*-free:

\begin{quote}
\small
\begin{verbatim}
function EVAL_BOOTSTRAP_VALUE(ckpt):
    for episode in test_episodes:
        obs = env.reset()
        goal_enc = jepa.encode(goal_observation)

        for step in 1 .. fixed_budget:
            for action in cardinal_actions:
                z_future[action] = jepa.rollout(obs, repeat(action, K))
                score[action] = V(z_future[action], goal_enc)
            action = argmax(score - revisit_penalty)
            obs = env.step(action)
            if reached_goal: success
\end{verbatim}
\end{quote}

\section{Pourquoi ce bootstrap est plus defendable}

\begin{center}
\begin{tabular}{p{0.28\linewidth}p{0.30\linewidth}p{0.30\linewidth}}
\toprule
Question & A* distille & Bootstrap WM \\
\midrule
Collecte de donnees & professeur A* optimal myopique & exploration locale sans solveur \\
Signal de cout & position A* \(N\) pas plus loin & reachability/value hindsight \\
Planification eval & A*-free & A*-free \\
Budget eval & proportionnel a A* initial & fixe, sans A* \\
Narration jury & bon controle appris & generalisation par modele du monde \\
\bottomrule
\end{tabular}
\end{center}

\section{Risques}

\begin{itemize}
  \item L'exploration peut ne pas atteindre assez souvent des etats proches du goal.
  \item La value hindsight apprend surtout la connectivite locale si les trajectoires
        ne couvrent pas de longs chemins.
  \item Le WM peut apprendre une dynamique locale correcte sans apprendre une strategie
        globale robuste.
  \item Le resultat doit donc etre compare a \code{probe\_pos} et \code{repr\_dist}
        sous le meme budget, pas seulement aux baselines A*.
\end{itemize}

\section{Etat de la run bootstrap}

Run bootstrap lancee sur Dalia:

\begin{quote}
\begin{tabular}{@{}ll@{}}
\toprule
Champ & Valeur \\
\midrule
Job Slurm & \code{75812} \\
Noeud & \code{dalianvl01} \\
Run dir & \code{$EBJEPA_WORK/runs/dynamic\_maze\_bootstrap\_20260620\_051011} \\
Config & \code{train\_dynamic\_maze\_bootstrap.yaml} \\
Collecte & \code{local\_goal\_frontier}, \code{layout\_use\_astar\_filter=false} \\
Budget train & \(5\times 16384 = 81920\) trajectoires nominales \\
Eval & \code{learned\_value}, \code{probe\_pos}, \code{repr\_dist}, 64 episodes chacun \\
\bottomrule
\end{tabular}
\end{quote}

Etat final observe le 20 juin 2026:

\begin{itemize}
  \item \done le job Slurm \code{75812} est \code{COMPLETED};
  \item \done duree totale: \code{01:32:25} sur un GPU \code{B200};
  \item \done le warm-up du pipeline stream est termine;
  \item \done la config est chargee et confirmee: \code{teacher\_policy=local\_goal\_frontier},
        \code{layout\_use\_astar\_filter=False}, \code{value\_mode=multi\_horizon\_td};
  \item \done checkpoints produits jusqu'a \code{e-4.pth.tar} et \code{latest.pth.tar};
  \item \done evaluations \code{learned\_value/probe\_pos/repr\_dist} produites;
  \item \done GIFs des premiers episodes copies dans \code{outputs/dynamic\_maze\_bootstrap\_20260620\_051011}.
\end{itemize}

\begin{center}
\begin{tabular}{lrrrr}
\toprule
Objectif & Success & Steps moyens & Blocked moves & Manhattan final \\
\midrule
\code{learned\_value} & 70.3\% & 289.5 & 88.6 & 7.0 \\
\code{probe\_pos} & 65.6\% & 323.2 & 84.7 & 6.8 \\
\code{repr\_dist} & 75.0\% & 244.9 & 85.1 & 4.1 \\
\bottomrule
\end{tabular}
\end{center}

Reference A* sequentielle sur la meme distribution non filtree:

\begin{center}
\begin{tabular}{lrrrrr}
\toprule
Politique & Success & Efficiency & Steps moyens & A* initial & Blocked \\
\midrule
\code{oracle\_replan\_seq\_rng} & 100.0\% & 0.963 & 67.5 & 73.1 & 0.0 \\
\bottomrule
\end{tabular}
\end{center}

Lecture en pourcentage du cas ideal A*: le \code{learned\_value} atteint
70.3\% du taux de succes ideal, mais seulement environ 23.3\% de l'efficacite
en nombre de deplacements. Le meilleur objectif actuel reste \code{repr\_dist}
avec 75.0\% de succes et environ 27.6\% d'efficacite en deplacements. La value
head apprend donc un signal utile, mais ne soutient pas encore completement la
claim Track 7 ``learned cost \(>\) latent distance''.

\section{Checklist}

\begin{itemize}
  \item \done Dynamic maze avec portes stochastiques et fog-of-war.
  \item \done Baselines A* strictes et GIFs.
  \item \done JEPA A* + sous-buts \(N=2,N=4\).
  \item \done Visualisations sans fog des portes et probabilites.
  \item \done Politique de collecte \code{local\_goal\_frontier} sans A*.
  \item \done Config bootstrap \code{train\_dynamic\_maze\_bootstrap.yaml}.
  \item \done Evaluateur bootstrap \code{eval\_bootstrap\_value.py}.
  \item \done Lancer le batch bootstrap sur Dalia.
  \item \done Suivre le job \code{75812} jusqu'au checkpoint WM/value.
  \item \done Comparer \code{learned\_value}, \code{probe\_pos}, \code{repr\_dist}.
  \item \done Ajouter GIFs et resultats bootstrap dans le depot.
  \item \todo Ameliorer \code{learned\_value} pour battre \code{repr\_dist}.
\end{itemize}

\end{document}
